\section{Two affine ends and one field parameter}
Put $L(q)=\mathrm{PSL}_2(q)$ and $P(q)=\mathrm{PGL}_2(q)$.
We retain the standard theorem that their automorphisms are induced by
projectivities and field automorphisms (with the usual small identifications).
No subgroup-classification theorem is needed in the following construction.
\begin{theorem}\label{E.PointedTriples}
Let $F=\mathbb F_q$, let $t\in F\setminus\{0,1,-1\}$ generate $F$ as a
field over its prime field, and put $k=|t|$. Consider the projective matrices
\[
 A=\begin{pmatrix}-t^{-1}&1\\0&1\end{pmatrix},\qquad
 B=\begin{pmatrix}t&0\\0&1\end{pmatrix},\qquad
 C=\begin{pmatrix}1&0\\1+t&-t\end{pmatrix}.
\]
Their pair products $AB,BC$ and triple product $ABC$ are nonidentity
involutions. Set $D=\langle-1,t\rangle\le F^\times$.
The groups $\langle A,B\rangle$ and $\langle B,C\rangle$ are the affine
groups $F^+\rtimes D$ fixing infinity and zero respectively; their
intersection is the diagonal group $D$. Consequently the rank-four
rotation intersection conditions hold exactly when $-1\in\langle t\rangle$.
In that case $\langle A,B,C\rangle$ equals $L(q)$ if $t$ is a square,
and $P(q)$ otherwise; in characteristic two these groups coincide.
The resulting polytope has type
\[
 (|{-t}|,k,|{-t}|).
\]
It is chiral exactly when no field automorphism sends $t$ to $t^{-1}$.
For a fixed admissible $k$, the number of isomorphism classes of the
ordered chiral markings is $\varphi(k)/[F:\mathbb F_p]$.
\end{theorem}
\begin{proof}
The displayed products have trace zero and are nonscalar; direct
multiplication proves the three involution relations in every characteristic.
The first end fixes infinity. It contains $B$ and
$R=AB=\left(\begin{smallmatrix}-1&1\\0&1\end{smallmatrix}\right)$.
The product $RBRB^{-1}$ is translation by $1-t$. Conjugation by $B$
scales translations by $t$, so its translation subgroup contains
$(1-t)\mathbb F_p[t]=F$. Its group of slopes is precisely $D$.
This proves the first affine description. After interchanging zero and
infinity, the second end contains a diagonal of slope $t^{-1}$ and an
affine map of slope $-t$ and translation $1+t$. The same commutator
argument supplies all translations, since $t\ne\pm1$, and proves its
description. An element in both ends fixes zero and infinity, so is
diagonal; both ends contain exactly the slopes $D$ there.

The nonidentity powers of $A$ fix precisely infinity and $t/(1+t)$;
those of $B$ fix infinity and zero. Thus their cyclic intersection is
trivial by the uniqueness of a projectivity on three distinct points.
Likewise $\langle B\rangle\cap\langle C\rangle=1$, since the fixed
points of $C$ are zero and one. Together with the computed end
intersection this proves precisely the rank-four intersection conditions.

The two ends contain the full opposite elementary unipotent groups.
These generate $L(q)$: elementary row operations reduce any determinant-one
matrix to a diagonal matrix, and such a diagonal is also a product of
upper and lower elementary matrices. The determinants of $A,B,C$ have
classes $[-t^{-1}],[t],[-t]$. Under the intersection hypothesis,
$-1\in\langle t\rangle$, their classes generate just $[t]$.
This proves the claimed whole group, without a maximal-subgroup test.
The eigenvalue ratios give the type. In odd characteristic the intersection
hypothesis says $k$ is even; then $|{-t}|=k$ for $4\mid k$ and
$|{-t}|=k/2$ for $k\equiv2\pmod4$. In characteristic two it is $k$.

For chirality use the retained rank-four rotation criterion: a reversing
automorphism would send $A$ to $A^{-1}$, $B$ to $A^2B$ and fix $C$.
Write it as a projectivity after a field automorphism $\sigma$.
Since $C$ is semisimple with eigenvalue ratio $-t$ and has order greater
than two, fixing its projective conjugacy class forces
$t^\sigma\in\{t,t^{-1}\}$. If $t^\sigma=t$, the full-field hypothesis
makes $\sigma$ the identity. A projectivity centralizing $C$ then fixes
zero and one individually, whereas one inverting $A$ must interchange
infinity and $u=t/(1+t)$. There is no such projectivity: writing it as
$\left(\begin{smallmatrix}a&0\\c&d\end{smallmatrix}\right)$, these
conditions give $a=c+d$, $a=uc$ and $d=-uc=-a$. Thus $c=2a$ and
$u=1/2$, forcing $t=1$; in characteristic two they already give $c=0$,
contradicting the interchange. This proves impossibility of inner reversal.

If instead $t^\sigma=t^{-1}$, then $\sigma^2=1$ because $t$ generates
$F$. The matrix $R$ above is defined over the prime field and satisfies
$R^2=I$. Direct substitution gives
\[
 RA^\sigma R=A^{-1},\qquad RB^\sigma R=A^2B,
 \qquad RC^\sigma R=-t^{-1}C.
\]
Thus $X\mapsto RX^\sigma R$ is the required involutory projective
automorphism, and the polytope is directly regular. This establishes
both directions of the chirality assertion.

Finally an isomorphism of two ordered markings sends the first end to
the first end and the second to the second. Their unique fixed points
are infinity and zero, so the projective part of its semilinear action
fixes these individually. It is diagonal, and comparison of $B$ shows
that the new parameter is $t^\sigma$. Conversely field automorphisms
send the whole displayed triple to the triple with parameter $t^\sigma$.
Every such orbit has length $[F:\mathbb F_p]$, and there are
$\varphi(k)$ elements of order $k$. This gives the stated number.
\end{proof}

\begin{proposition}\label{E.PointedCompleteness}
Suppose a rank-four rotation marking generates $L(q)$ or $P(q)$ and each
of its two ends fixes a point of the projective line. Then, up to
projective conjugation, it has the displayed form in
Theorem~\ref{E.PointedTriples}, with $t$ generating the whole field.
\end{proposition}
\begin{proof}
The two fixed points differ, since the whole group fixes no point.
Their common generator $B$ fixes both, so conjugate them to infinity
and zero and write $B=\operatorname{diag}(t,1)$.
The first and last generators have the forms
$A=\left(\begin{smallmatrix}\lambda&\mu\\0&1\end{smallmatrix}\right)$
and $C=\left(\begin{smallmatrix}1&0\\\nu&\kappa\end{smallmatrix}\right)$.
The involution relations give $\lambda=-t^{-1}$, $\kappa=-t$ and
$\mu\nu=1+t$. Here $t\ne1$, since $B$ is nonidentity, and $t\ne-1$:
otherwise $\mu\nu=0$ and the whole group fixes zero or infinity.
Thus $\mu\ne0$, and conjugation by $\operatorname{diag}(1,\mu)$
normalizes $\mu$ to one. This gives the displayed matrices.
They all lie in $P(\mathbb F_p(t))$, so the group-order formulas force
$\mathbb F_p(t)=F$. (For a proper subfield of size $s$, $q\ge s^2$
and $|L(q)|>|P(s)|$.) The preceding theorem applies.
\end{proof}

The full-field condition is $k\nmid p^e-1$ for $0<e<d$, where
$q=p^d$; the absence of field reversal is $k\nmid p^e+1$ for the same
range. Thus this one theorem gives both source pointed-parabolic
classifications, including their counts. For $P(q)$ with odd $q$, add
$k\nmid(q-1)/2$; for $L(q)$ with odd $q$, add $k\mid(q-1)/2$ and
$k$ even. The latter implies $q\equiv1\pmod4$. In characteristic two
there is no determinant-square distinction and no even-$k$ requirement.

There is a proof-route correction to the source's subfield sentence.
When $t$ generates $F$ but a field automorphism inverts it, the group
can still be all of $L(q)$ or $P(q)$. The exclusion from the chiral list
comes from the semilinear reversal just constructed. For example, an
element $t\in\mathbb F_{25}$ of order six gives $L(25)$, of order
$7800$, with affine ends of order $150$ and intersection of order six;
Frobenius inversion makes this marking regular. These orders follow
uniformly from the theorem, not from a sample extrapolation.
